
\begin{sublemme}
\label{XIII.3.1.2}
Sean $S$ un esquema localmente noetheriano, $f\colon X\to S$ un
morfismo, $F\to G$ un monomorfismo de haces de grupos sobre~$X$,
$\Psi$ (\resp $\Psi_1$) el stack de los torsores bajo~$F$ (\resp bajo~$G$),
$\Phi=f_*\Psi$, $\Phi_1=f_*\Psi_1$. Se supone dado un haz sobre~$S$
representable por un $S$-esquema étale de tipo finito~$T$ y un morfismo
sobreyectivo
$$
a\colon T\to S\Phi_1\simeq\R^1f_*G,
$$
de modo que exista un torsor~$Q$ sobre~$X_T=X\times_ST$, bajo~$G|X_T$,
cuya clase en~$\R^1f_*G(T)$ sea la imagen por~$a$ de la sección identidad
de~$T(T)$. Sea $f_T\colon X_T\to T$ el morfismo canónico y pongamos
$F_T=F|X_T$. Se supone que $\Phi_1$ es $1$-constructible y que
$f_{T*}(Q/F_T)$ es constructible. Entonces $\Phi$ es $1$-constructible.
\end{sublemme}

En efecto, basta repetir la demostración de~\Ref{XIII.3.1.1},
sustituyendo la existencia de los morfismos~$\Psi_i$ por la de los
isomorfismos
$$
\ifthenelse{\boolean{orig}}
{F'_i\isomto f_*(Q/F_T)|S'.}
{F'_i\isomto f_{T*}(Q/F_T)|S'.}
$$
